\documentclass[english]{article}
\usepackage[T1]{fontenc}
\usepackage[utf8]{inputenc}
\setcounter{secnumdepth}{-2}
\setlength{\parindent}{0bp}
\usepackage{amsmath}
\usepackage{amsthm}
\usepackage{amssymb}

\makeatletter
\theoremstyle{definition}
\newtheorem*{problem*}{\protect\problemname}

\@ifundefined{date}{}{\date{}}
\usepackage{graphicx}
\usepackage{geometry}
\newcommand{\limi}{\underset{n\rightarrow\infty}{\lim}}
\newcommand{\limxz}{\underset{x\rightarrow x_{0}}{\lim}}
\newcommand{\limfxz}{\underset{x\rightarrow x_{0}}{\lim}f(x)}
\newcommand{\limfxm}{\underset{x\rightarrow x_{0}^{-}}{\lim}f(x)}
\newcommand{\limfxp}{\underset{x\rightarrow x_{0}^{+}}{\lim}f(x)}
\newcommand{\limxm}{\underset{x\rightarrow x_{0}^{-}}{\lim}}
\newcommand{\limxp}{\underset{x\rightarrow x_{0}^{+}}{\lim}}
\newcommand{\sqrm}{M_{m\times m}(\mathbb{F})}
\newcommand{\seqa}{(a_{n})_{n=1}^{\infty}}
\newcommand{\seqx}{(x_{n})_{n=1}^{\infty}}
\newcommand{\funcdr}{f:D\rightarrow\mathbb{R}}
\newcommand{\der}{\limxz\frac{f(x)-f(x_{0})}{x-x_{0}}}
\usepackage[all]{pdfprivacy}

\makeatother

\usepackage{babel}
\providecommand{\problemname}{Problem}

\begin{document}
\title{{\Large\vspace{-6em}Semester 2026A, Quiz 4\vspace{-3em}}}

\maketitle
\rule[0.5ex]{1\columnwidth}{1pt}
\begin{abstract}
Quiz 4 in the course Infinitesimal Calculus (1), semester 2026 A.

Quiz instructions: The duration of the quiz is 36 minutes.

The use of calculators and of any other auxiliary material is forbidden.
\\

This is a quiz that I translated and solved.

The original quiz was written by the Infi Calculus 1 course lecturers
at HUJI.
\end{abstract}
\medskip{}

\rule[0.5ex]{1\columnwidth}{1pt}

\medskip{}

\begin{problem*}[Question 1 - 65 points]
 Prove or disprove:

If $f:\mathbb{R}\to\mathbb{R}$ is a continuous function at $x_{0}\in\mathbb{R}$
and in every punctured neighborhood of $x_{0}$ there exists a point
$z\in\mathbb{R}$ s.t. $f(z)=z^{2}$, then $f(x_{0})=x_{0}^{2}$.
\end{problem*}
\begin{proof}
\textbf{using the $\varepsilon-\delta$ definition:}\medskip{}

1. Since $f$ is continuous at $x_{0}$, 
\[
\forall\varepsilon>0\exists\delta>0\ \text{s.t.}\ \forall x\in\left(x_{0}-\delta,x_{0}+\delta\right):\left|f(x)-f(x_{0})\right|<\varepsilon
\]

2. We shall define $g:\mathbb{R}\rightarrow\mathbb{R}$ by $g(x)=x^{2}$. 

The function $g$ is continuous on $\mathbb{R}$ because it is a polynomial,
hence the same continuity characterisation as in point 1 applies to
it.

3. Let us assume for the sake of contradiction that $f(x_{0})\ne x_{0}^{2}$.
Immediately, $f(x_{0})\ne g(x_{0})$.

4. The law of trichotomy establishes that either $f(x_{0})>g(x_{0})$
or $f(x_{0})<g(x_{0})$.

We will assume without loss of generality that $f(x_{0})>g(x_{0})$.

Hence, $f(x_{0})-g(x_{0})>0$.

5. The previous point allows us to apply the known technique of halving
the distance to separate these values and set $\varepsilon:=\frac{f(x_{0})-g(x_{0})}{2}>0$
in the property of points 1 and 2, which gives us corresponding deltas
$\delta_{1},\delta_{2}>0$.

We shall define $\delta_{3}:=\min\left\{ \delta_{1},\delta_{2}\right\} $
and get the combined property: 
\[
\forall x\in\left(x_{0}-\delta_{3},x_{0}+\delta_{3}\right)\setminus\left\{ x_{0}\right\} :\left|f(x)-f(x_{0})\right|<\frac{f(x_{0})-g(x_{0})}{2}\ \wedge\left|g(x)-g(x_{0})\right|<\frac{f(x_{0})-g(x_{0})}{2}
\]

Expanding the absolute values gives:
\[
\frac{g(x_{0})-f(x_{0})}{2}<f(x)-f(x_{0})<\frac{f(x_{0})-g(x_{0})}{2}\ \wedge\ \frac{g(x_{0})-f(x_{0})}{2}<g(x)-g(x_{0})<\frac{f(x_{0})-g(x_{0})}{2}
\]

Isolating $f(x)$ and $g(x)$ on the relevant sides provides a lower
bound for $f(x)$ and an upper bound for $g(x)$:
\[
\frac{g(x_{0})+f(x_{0})}{2}<f(x)\quad\wedge\quad g(x)<\frac{f(x_{0})+g(x_{0})}{2}
\]

Thus, in particular, 
\[
\forall x\in\left(x_{0}-\delta_{3},x_{0}+\delta_{3}\right)\setminus\left\{ x_{0}\right\} :g(x)<f(x)
\]

6. This means that 
\[
\forall x\in\left(x_{0}-\delta_{3},x_{0}+\delta_{3}\right)\setminus\left\{ x_{0}\right\} :f(x)\ne x^{2}
\]
In other words, $\left(x_{0}-\delta_{3},x_{0}+\delta_{3}\right)\setminus\left\{ x_{0}\right\} $
is a punctured neighborhood of $x_{0}$ which contradicts the assumption
that every punctured neighborhood of $x_{0}$ contains a point $z\in\mathbb{R}$
s.t.$f(z)=z^{2}$.
\end{proof}
\newpage{}
\begin{proof}
\textbf{using Heine's sequences:}\medskip{}

1. By the assumption, for each $n\in\mathbb{N}$, there exists a point
\[
x_{n}\in\left(x_{0}-\frac{1}{n},x_{0}+\frac{1}{n}\right)\setminus\left\{ x_{0}\right\} \quad\text{s.t.\ensuremath{\quad}}f(x_{n})=x_{n}^{2}
\]
Thus, we have a sequence $\left(x_{n}\right)_{n=1}^{\infty}$ such
that for all $n\in\mathbb{N},$ 
\[
x_{0}-\frac{1}{n}<x_{n}<x_{0}+\frac{1}{n}\ \wedge\ x_{n}\ne x_{0}\ \wedge\ f(x_{n})=x_{n}^{2}
\]

2. Since $\underset{n\rightarrow\infty}{\lim}\left(x_{0}-\frac{1}{n}\right)=\underset{n\rightarrow\infty}{\lim}\left(x_{0}+\frac{1}{n}\right)=x_{0}$,
the squeeze theorem implies $\underset{n\rightarrow\infty}{\lim}x_{n}=x_{0}$.

3. Points 1 and 2 show that $\left(x_{n}\right)_{n=1}^{\infty}$ satisfies
the conditions of a Heine sequence for continuity.

By the assumption, $f$ is continuous at $x_{0}$, therefore $\underset{n\rightarrow\infty}{\lim}f(x_{n})=f(x_{0})$.

4. According to the arithmetic of sequence limits,
\[
\underset{n\rightarrow\infty}{\lim}f(x_{n})=\underset{n\rightarrow\infty}{\lim}x_{n}^{2}=\underset{n\rightarrow\infty}{\lim}x_{n}\cdot\underset{n\rightarrow\infty}{\lim}x_{n}=x_{0}\cdot x_{0}=x_{0}^{2}
\]
 5. Therefore, $f(x_{0})=x_{0}^{2}$, as required. 
\end{proof}
\medskip{}

\begin{problem*}[Question 2 - 35 points]
 Let $(a_{n})_{n=1}^{\infty}$ be a sequence of natural numbers that
satisfies $\limi\frac{a_{n}}{n}=\infty$.
\end{problem*}
Prove that $\limi\left(1+\frac{1}{n}\right)^{a_{n}}=\infty$.
\begin{proof}
Since $\limi\frac{a_{n}}{n}=\infty$, the arithmetic of sequence limits
ensures $\limi\left(1+\frac{a_{n}}{n}\right)=\infty$.

Additionally, because $a_{n}\in\mathbb{N}$ and $\frac{1}{n}>-1$
for all $n\in\mathbb{N}$, Bernoulli's inequality holds and 
\[
\forall n\in\mathbb{N}:\left(1+\frac{1}{n}\right)^{a_{n}}\ge1+\frac{a_{n}}{n}
\]

Recall the slice theorem: for two sequences $(b_{n})_{n=1}^{\infty}$
and $(c_{n})_{n=1}^{\infty}$, if there exists $N\in\mathbb{N}$ such
that for every $n>N$ the inequality $b_{n}\ge c_{n}$ holds, and
$\limi c_{n}=\infty$, then $\limi b_{n}=\infty$. 

Therefore, $\limi\left(1+\frac{1}{n}\right)^{a_{n}}=\infty$ as required. 
\end{proof}

\end{document}
